\section*{Розділ I. Загальні положення}
\addcontentsline{toc}{section}{Розділ I. Загальні положення}

\subsection*{1.1. Мета та сфера застосування}
\addcontentsline{toc}{subsection}{1.1. Мета та сфера застосування}
    1.1.1. Цей Тимчасовий виборчий регламент (далі -- Регламент) визначає порядок організації та проведення перших виборів до органів студентського самоврядування Державного некомерційного підприємства ``Державний університет ``Київський авіаційний інститут'' (далі -- Університет) після набуття чинності Положенням про органи студентського самоврядування.

    1.1.2. Мета цього Регламенту -- вирішення проблеми циклічної залежності між необхідністю проведення виборів для формування органів студентського самоврядування та потребою у функціонуючих органах для організації цих виборів.

    1.1.3. Регламент застосовується виключно для проведення перших виборів делегатів до Конференції студентів Університету (КСУ).

\subsection*{1.2. Принципи тимчасового регламенту}
\addcontentsline{toc}{subsection}{1.2. Принципи тимчасового регламенту}
    1.2.1. Діяльність за цим Регламентом ґрунтується на принципах: законності, демократичності, прозорості, рівності виборчих прав студентів, альтернативності, добровільності участі у виборах та колегіальності прийняття рішень.

    1.2.2. Цей Регламент забезпечує безперервність та наступність у формуванні органів студентського самоврядування відповідно до Положення про органи студентського самоврядування Університету.

    1.2.3. Процедури, визначені цим Регламентом, мають тимчасовий характер і припиняють свою дію після досягнення мети -- формування повноцінної системи органів студентського самоврядування.

\subsection*{1.3. Правова основа та співвідношення з іншими нормативними актами}
\addcontentsline{toc}{subsection}{1.3. Правова основа та співвідношення з іншими нормативними актами}
    1.3.1. Правову основу цього Регламенту становлять Конституція України, Закон України ``Про вищу освіту'', Статут Університету, Положення про органи студентського самоврядування Університету та інші нормативно-правові акти України.

    1.3.2. У питаннях проведення перших виборів до органів студентського самоврядування цей Регламент має спеціальний характер та застосовується з пріоритетом над відповідними нормами Положення про органи студентського самоврядування, Положення про Центральну виборчу комісію студентів та інших нормативних актів студентського самоврядування.

    1.3.3. Питання, не врегульовані цим Регламентом, вирішуються відповідно до Положення про органи студентського самоврядування Університету, Статуту Університету та загальних принципів права.

    1.3.4. У разі виникнення суперечностей між нормами цього Регламенту та Положенням про органи студентського самоврядування щодо процедур перших виборів, застосовуються норми цього Регламенту.

\subsection*{1.4. Строк дії та умови припинення}
\addcontentsline{toc}{subsection}{1.4. Строк дії та умови припинення}
    1.4.1. Цей Регламент діє з моменту його затвердження Загальними зборами студентів Університету до проведення перших виборів делегатів КСУ та скликання першого засідання КСУ.

    1.4.2. Регламент припиняє свою дію після виконання наступної умови: проведення перших виборів делегатів КСУ та скликання першого засідання КСУ.

    1.4.3. Незалежно від виконання умов пункту 1.4.2, дія цього Регламенту автоматично припиняється через 6 (шість) місяців з дня його затвердження.

    1.4.4. Про припинення дії цього Регламенту Тимчасова виборча комісія повідомляє студентську спільноту через офіційні інформаційні канали не пізніше ніж за 3 дні до припинення дії.

\subsection*{1.5. Основні поняття}
\addcontentsline{toc}{subsection}{1.5. Основні поняття}
    1.5.1. У цьому Регламенті терміни вживаються у значеннях, визначених Положенням про органи студентського самоврядування Університету, а також:

        \begin{enumerate}[label=\alph*)]
            \item \textbf{Тимчасова виборча комісія (ТВК)} -- колегіальний орган, створений для організації та проведення перших виборів делегатів КСУ відповідно до цього Регламенту;
            \item \textbf{Перші вибори} -- вибори делегатів КСУ, що проводяться вперше після набуття чинності Положенням про органи студентського самоврядування;
            \item \textbf{Загальні збори студентів (ЗЗС)} -- зібрання студентів Університету, скликане для затвердження цього Регламенту та формування ТВК;
            \item \textbf{Виборчий округ} -- територіальна або структурна одиниця Університету, в межах якої проводяться вибори делегатів КСУ.
        \end{enumerate}